% Осень26!A18; нумерация внутри столбца без пустых строк.
% Источник: лекция 6, §7 (строка 389)
% Источник: лекция 7, §9 (строка 272)
\ticket{16}{Теорема о гладкости обратного отображения}
\label{ticket:16}

\begin{theorem}
Пусть $O\subset\mathbb R^m$ открыто, $T:O\to\mathbb R^m$ принадлежит классу $C^r$, $r\in\mathbb N\cup\{\infty\}$, биективно на $T(O)$ и $\det T'(x)\ne0$ на $O$.
Тогда $S=T^{-1}\in C^r(T(O))$ и
\[
  S'(y)=\bigl(T'(S(y))\bigr)^{-1}.
\]
\end{theorem}

\begin{proof}
По \hyperref[ticket:14]{теореме о сохранении области} $T$ переводит открытые множества в открытые.
Поэтому $T(O)$ открыто и $S$ непрерывно.
Пусть $y_0=T(a)$, $A=T'(a)$.
При $x\to a$
\[
  T(x)-T(a)=A(x-a)+r(x),\qquad r(x)=o(|x-a|).
\]
По \hyperref[ticket:13]{лемме об оценке приращения} вблизи $a$ имеем $|x-a|\le C|T(x)-T(a)|$.
Подставляя $x=S(y)$, получаем
\[
  S(y)-a=A^{-1}(y-y_0)-A^{-1}r(S(y)).
\]
Непрерывность $S$ и указанная оценка дают $r(S(y))=o(|y-y_0|)$.
Следовательно, $S$ дифференцируемо и $S'(y_0)=A^{-1}$.
В формуле $S'=(T'\circ S)^{-1}$ все операции непрерывны, включая \hyperref[ticket:08]{обращение матрицы}; значит, $S\in C^1$.

Пусть теперь $T\in C^{k+1}$ и уже доказано $S\in C^k$.
Элементы матрицы $T'\circ S$ принадлежат $C^k$.
По формулам Крамера элементы обратной матрицы --- отношения многочленов от этих элементов с ненулевым знаменателем $\det(T'\circ S)$.
Поэтому $S'\in C^k$, то есть $S\in C^{k+1}$.
Индукция доказывает утверждение для конечных $r$, а случай $C^\infty$ следует из всех конечных порядков.
\end{proof}
